\chapter{Pricers de Monte-Carlo en pratique}\label{ch:dv:monte-carlo-pricers-in-practice}

Une option bermudéenne sur le meilleur de cinq actions n'a aucune grille aux différences finies qui tienne en mémoire~: cent points par
dimension, soit dix milliards de nœuds. Un moteur de Monte-Carlo la price en quelques secondes, mais il se trompe par construction
sur la décision d'exercice. Il estime la valeur de l'attente par une régression sur les trajectoires simulées, ce qui n'est
qu'une approximation de la vraie valeur de continuation, si bien que la politique d'exercice suivie est sous-optimale et que son prix
est une borne inférieure. Un bon desk publie donc deux nombres, une borne inférieure issue de la politique et une borne supérieure issue d'une méthode duale, et l'écart entre
les deux mesure ce que la régression laisse encore passer. Ce chapitre construit l'outillage qui entoure la méthode
de Monte-Carlo du livre~4~: simulation des modèles, exercice anticipé par régression et sa borne duale, trois façons de calculer les \omterm{def:dv:greeks-and-the-hedging-pnl:greeks}{grecques},
nombres à discrépance faible avec le pont brownien, et l'ingénierie qui en fait un pricer.

\section{Simulation des modèles}

Les actifs lognormaux se simulent exactement, date après date. Tout autre modèle est discrétisé, et l'erreur du schéma est un
biais qui s'ajoute à l'erreur statistique et ne s'annule pas en moyennant. La variance de Heston est le cas difficile classique. Un pas d'Euler
peut rendre la variance négative, et la tronquer à zéro laisse un biais qui ne décroît que lentement avec le pas de temps. Andersen
a conçu un schéma qui reproduit les deux premiers moments de la transition exacte de la variance.

\begin{definition}[Schéma quadratique-exponentiel]\label{def:dv:monte-carlo-pricers-in-practice:qe}
Le \emph{schéma quadratique-exponentiel}\index{schéma quadratique-exponentiel} (quadratic-exponential scheme, QE) simule la variance de Heston sur un pas en égalant
la moyenne $m$ et la variance $s^2$ de sa loi conditionnelle exacte~: lorsque $\psi=s^2/m^2$ est petit, il tire $a(b+Z)^2$, un
carré non central, mis à l'échelle, d'une loi normale~; lorsque $\psi$ est grand, un mélange d'une masse en zéro et d'une loi exponentielle. Le log-prix utilise alors la
variance intégrée implicite des variances initiale et finale.
\end{definition}

\begin{example}[Euler contre QE]\label{ex:dv:monte-carlo-pricers-in-practice:qe}
Une option d'achat à la monnaie d'un an dans le \omterm{def:dv:stochastic-volatility:heston}{modèle de Heston} de base du chapitre 10 ($v_0=\bar v=0.04$, $\kappa=1.5$, $\eta=0.6$, $\rho=-0.7$) vaut 6.751 par
inversion de Fourier. Avec 400\,000 trajectoires et un Euler à troncature complète, le biais vaut 0.857 avec 4 pas par an, 0.373 avec 8, 0.140 avec 16 et
0.041 avec 32. Avec QE, le prix est à moins d'une erreur type (0.014) de la valeur exacte pour tous les nombres de pas, dès 4 pas par an
(\cref{fig:dv:monte-carlo-pricers-in-practice:heston}). QE permet de diviser par dix au moins le nombre de pas, donc le temps de calcul.
\end{example}

\begin{omfigure}
\begin{tikzpicture}
\begin{axis}[width=10.5cm,height=5cm,xmode=log,ymode=log,log basis x=2,
  xlabel={pas de temps par an},ylabel={biais absolu},
  tick label style={font=\small},label style={font=\small},
  xmin=3,xmax=40,ymin=1e-4,ymax=2,grid=major,grid style={gray!25},
  xtick={4,8,16,32},xticklabels={4,8,16,32},
  legend columns=3,legend cell align=left,legend style={font=\footnotesize,at={(0.5,-0.3)},anchor=north,draw=none,fill=none,/tikz/every even column/.append style={column sep=8pt}},
  table/col sep=comma]
\addplot[omProp,very thick,mark=*,mark size=1.6pt] table[x=steps,y=euler]{figdata/derivatives/23-monte-carlo-pricers-in-practice/heston_bias.csv};
\addplot[omThm,very thick,mark=square*,mark size=1.6pt] table[x=steps,y=qe]{figdata/derivatives/23-monte-carlo-pricers-in-practice/heston_bias.csv};
\addplot[gray,thick,dashed] table[x=steps,y=se]{figdata/derivatives/23-monte-carlo-pricers-in-practice/heston_bias.csv};
\legend{Euler à troncature complète,schéma quadratique-exponentiel,erreur type}
\end{axis}
\end{tikzpicture}

\omcaption{Biais d'un call à la monnaie d'un an dans un \omterm{def:dv:stochastic-volatility:heston}{modèle de Heston} dont la variance peut atteindre zéro, par pas de temps, par rapport au prix de Fourier
(400\,000 trajectoires). Le biais d'Euler décroît lentement~; celui de QE reste dans l'erreur statistique à chaque nombre de pas. Données~: le
tutoriel.}\label{fig:dv:monte-carlo-pricers-in-practice:heston}
\end{omfigure}

\section{Exercice anticipé par régression}

Le chapitre 6 évaluait à rebours les options bermudéennes et américaines sur un arbre, où la valeur de continuation à chaque nœud est une moyenne sur ses
nœuds suivants. Une simulation avance dans le temps et ne dispose d'aucune moyenne de ce type. Longstaff et Schwartz l'ont estimée par régression.

\begin{definition}[Méthode de Longstaff--Schwartz]\label{def:dv:monte-carlo-pricers-in-practice:lsm}
La \emph{méthode de Longstaff--Schwartz}\index{méthode de Longstaff--Schwartz} (Longstaff--Schwartz method) évalue une option bermudéenne par simulation~: en remontant les
dates d'exercice, elle régresse les flux actualisés que chaque trajectoire réalise à partir de la date suivante, sur un jeu de fonctions de base de la
variable d'état, en ne retenant que les trajectoires dans la monnaie, et exerce lorsque la valeur d'exercice dépasse la valeur de continuation ajustée~; la politique est alors
appliquée à des trajectoires indépendantes pour fournir une borne inférieure.
\end{definition}

Le prix issu d'une politique estimée est biaisé à la baisse, car aucune politique ne fait mieux que la politique optimale. Ajuster et évaluer sur les mêmes trajectoires ajoute un
biais à la hausse, car la régression a déjà vu les trajectoires qu'elle évalue ensuite. La construction ajuste sur un premier jeu de trajectoires et évalue sur un autre, ce qui garde au résultat son
caractère de borne inférieure. Son niveau dépend de la base retenue. Une régression incapable de représenter la valeur de continuation conduit à exercer au
mauvais moments.

\begin{definition}[Borne supérieure duale]\label{def:dv:monte-carlo-pricers-in-practice:dual}
La \emph{borne supérieure duale}\index{borne supérieure duale} (dual upper bound) d'une option bermudéenne est $\E\bigl[\max_j(h_j-M_j)\bigr]$ pour une martingale $M$ avec $M_0=0$
et des valeurs d'exercice actualisées $h_j$~; elle dépasse le prix pour toute martingale et lui est égale pour la partie martingale
du processus de valeur. L'algorithme d'Andersen--Broadie construit $M$ à partir d'une politique de borne inférieure par simulation emboîtée, de sorte que les
bornes inférieure et supérieure donnent ensemble un intervalle pour le prix.
\end{definition}

Dans la construction, la valeur de continuation de la politique $Q_j$ est estimée à chaque nœud externe par des trajectoires internes qui suivent la politique. Avec $L_j$ la
valeur d'exercice actualisée là où la politique exerce et $Q_j$ là où elle continue, les accroissements de la martingale valent $L_j-Q_{j-1}$. Les
estimations internes sont bruitées, et le maximum dans la borne transforme le bruit en biais à la hausse. La borne est valide pour un nombre quelconque de trajectoires internes, mais serrée seulement
pour un grand nombre.

\begin{example}[Le biais d'un nombre trop faible de trajectoires internes]\label{ex:dv:monte-carlo-pricers-in-practice:inner}
Un put bermudéen d'un an à dix dates d'exercice (spot et strike 100, 5~\%, 20~\%) vaut 6.0326 sur l'arbre du chapitre 6. La \omterm{def:dv:monte-carlo-pricers-in-practice:lsm}{méthode de Longstaff--Schwartz} avec une
base cubique donne 6.0325 (erreur type 0.023). La borne duale sur 1\,000 trajectoires externes vaut 6.134 avec 100 trajectoires internes par nœud, 6.054 avec 400 et
6.037 avec 1\,600 (\cref{fig:dv:monte-carlo-pricers-in-practice:bounds}, à droite).
\end{example}

\begin{example}[Un max-call bermudéen à trois actifs]\label{ex:dv:monte-carlo-pricers-in-practice:maxcall}
Un call sur le meilleur de trois actions indépendantes (chacune à 100, strike 100, volatilité 20~\%, rendement de dividende 10~\%, taux 5~\%), exerçable à neuf dates
sur trois ans. La base de la régression est un ensemble de polynômes en les deux plus grands prix, leur produit, le produit des trois et la
valeur intrinsèque. Avec cette base, Longstaff--Schwartz sur 200\,000 trajectoires fraîches donne 18.649 (erreur type 0.039), et la borne duale (1\,000 trajectoires externes,
1\,000 trajectoires internes) donne 18.713 (0.019). L'écart vaut 0.064, soit 0.3~\% du prix. Avec une base réduite au plus grand prix et à son carré, la borne inférieure
tombe à 18.137 et la borne supérieure monte à 18.805, soit un écart de 0.668, dix fois plus large
(\cref{fig:dv:monte-carlo-pricers-in-practice:bounds}, à gauche).
\end{example}

\begin{omfigure}
\begin{tikzpicture}
\begin{groupplot}[group style={group size=2 by 1,horizontal sep=1.6cm},
  width=6.6cm,height=5cm,
  tick label style={font=\small},label style={font=\small},grid=major,grid style={gray!25},table/col sep=comma]
\nextgroupplot[xlabel={base de régression},ylabel={prix},xmin=-0.6,xmax=1.6,ymin=18.0,ymax=18.95,
  xtick={0,1},xticklabels={riche,réduite},title={\small bornes du max-call},
  legend columns=1,legend cell align=left,legend style={font=\footnotesize,at={(0.5,-0.32)},anchor=north,draw=none,fill=none}]
\addplot[omDef,only marks,mark=square*,mark size=2.4pt,error bars/.cd,y dir=both,y explicit] table[x=i,y=lower,y error=lower_se]{figdata/derivatives/23-monte-carlo-pricers-in-practice/bounds.csv};
\addplot[omThm,only marks,mark=triangle*,mark size=2.8pt,error bars/.cd,y dir=both,y explicit] table[x=i,y=upper,y error=upper_se]{figdata/derivatives/23-monte-carlo-pricers-in-practice/bounds.csv};
\legend{Longstaff--Schwartz (inférieure),duale (supérieure)}
\nextgroupplot[xmode=log,log basis x=2,xlabel={trajectoires internes par nœud},ylabel={borne supérieure duale},xmin=70,xmax=2200,ymin=6.0,ymax=6.16,
  xtick={100,400,1600},xticklabels={100,400,1\,600},
  yticklabel style={/pgf/number format/fixed,/pgf/number format/precision=2},title={\small put bermudéen},
  legend columns=1,legend cell align=left,legend style={font=\footnotesize,at={(0.5,-0.32)},anchor=north,draw=none,fill=none}]
\addplot[omThm,very thick,mark=*,mark size=1.6pt] table[x=inner,y=upper]{figdata/derivatives/23-monte-carlo-pricers-in-practice/inner_bias.csv};
\addplot[black,dashed] table[x=inner,y=tree]{figdata/derivatives/23-monte-carlo-pricers-in-practice/inner_bias.csv};
\legend{borne duale,valeur sur l'arbre}
\end{groupplot}
\end{tikzpicture}

\omcaption{À gauche~: bornes inférieure et supérieure (barres d'une erreur type) pour un max-call bermudéen à trois actifs, avec une base de régression riche et une base de régression réduite~:
l'écart est le prix d'une base pauvre. À droite~: la borne duale d'un put bermudéen à un actif, en fonction du nombre de trajectoires internes par nœud, qui se rapproche
de la valeur sur l'arbre par le haut. Données~: le tutoriel.}\label{fig:dv:monte-carlo-pricers-in-practice:bounds}
\end{omfigure}

\section{Grecques~: trajectorielle, par rapport de vraisemblance, adjointe}

Le \omterm{def:dv:greeks-and-the-hedging-pnl:bump}{bump-and-reprice} d'un prix Monte-Carlo fournit des \omterm{def:dv:greeks-and-the-hedging-pnl:greeks}{grecques} dont la variance est celle d'une différence de deux nombres bruités. Les nombres aléatoires communs rendent la
différence beaucoup moins bruitée. Deux méthodes directes, dues à Broadie et Glasserman, évitent entièrement le bump et donnent des \omterm{def:dv:greeks-and-the-hedging-pnl:greeks}{grecques} à partir des mêmes trajectoires que le
prix.

\begin{definition}[Grecque trajectorielle]\label{def:dv:monte-carlo-pricers-in-practice:pathwise}
Une \emph{grecque trajectorielle}\index{grecque trajectorielle} dérive le paiement actualisé le long de chaque trajectoire simulée par rapport au paramètre, en gardant
les nombres aléatoires fixés, puis en calcule la moyenne~; elle exige que le paiement soit continu en fonction du paramètre (presque sûrement dérivable), et échoue pour
les digitales et les options à barrière.
\end{definition}

\begin{definition}[Grecque par rapport de vraisemblance]\label{def:dv:monte-carlo-pricers-in-practice:lr}
Une \emph{grecque par rapport de vraisemblance}\index{grecque par rapport de vraisemblance} dérive la densité de l'état simulé au lieu du paiement~: elle moyenne le
paiement actualisé multiplié par le score, c'est-à-dire la dérivée de la log-densité par rapport au paramètre~; elle fonctionne pour les paiements discontinus, au prix
d'une variance plus élevée.
\end{definition}

Le tableau compare les deux méthodes sur des options à la monnaie d'un an (cours au comptant et strike 100, taux 5~\%, volatilité 20~\%) avec 200\,000 trajectoires, erreurs types
entre crochets. La dernière colonne applique la méthode trajectorielle à la digitale lissée en un call spread.

\begin{center}\small
\begin{tabular}{lrrrr}
\toprule
un an & exact & trajectorielle & rapport de vraisemblance & lissée\\
\midrule
delta du call & 0.6368 & 0.6384 (0.0013) & 0.6413 (0.0033) & ---\\
\omterm{def:dv:greeks-and-the-hedging-pnl:greeks}{vega} du call & 37.52 & 37.72 (0.17) & 38.24 (0.62) & ---\\
delta de la digitale & 0.01876 & 0 (faux) & 0.01883 (0.00006) & 0.01895 (0.00021)\\
\omterm{def:dv:greeks-and-the-hedging-pnl:greeks}{vega} de la digitale & $-0.657$ & 0 (faux) & $-0.648$ (0.010) & $-0.663$ (0.007)\\
\bottomrule
\end{tabular}
\end{center}

Pour le call, la variance de l'estimateur trajectoriel est 6.5 fois plus faible pour le delta et 13 fois plus faible pour le \omterm{def:dv:greeks-and-the-hedging-pnl:greeks}{vega}. Pour la digitale, il renvoie zéro, puisque la
dérivée du paiement est nulle presque partout. Le rapport de vraisemblance donne le bon résultat. Le lissage de la digitale en un call spread de demi-largeur un rétablit
une estimation trajectorielle, au prix d'un biais propre au call spread.

La méthode trajectorielle dispose d'une implémentation efficace, qui compte surtout pour les grands portefeuilles. La méthode adjointe de Giles et Glasserman
parcourt en sens inverse, le long de chaque trajectoire, le même calcul trajectoriel~: elle donne les mêmes valeurs, et son avantage croît avec le nombre
d'entrées, comme les points d'une \omterm{def:dv:implied-volatility-and-its-surface:surface}{surface de volatilité}.

\begin{example}[Douze vegas pour le prix de trois passes]\label{ex:dv:monte-carlo-pricers-in-practice:adjoint}
Une \omterm{def:dv:asians-lookbacks-cliquets-and-forward-starts:asian}{option asiatique} arithmétique (un an, 48 pas, strike 100, 20~\%) est évaluée avec sa sensibilité à la volatilité dans chacun des douze compartiments mensuels.
Le mode direct demanderait une passe par compartiment. La passe adjointe (inverse) du livre~4, chapitre~28, parcourt une fois en sens inverse chaque trajectoire stockée
et accumule la dérivée du paiement par rapport à l'état et à la volatilité de chaque pas. Elle restitue les douze \omterm{def:dv:greeks-and-the-hedging-pnl:greeks}{vegas} et le delta pour un coût de
quelques passes directes. Le prix vaut 5.901 et le delta 0.593. Le \omterm{def:dv:greeks-and-the-hedging-pnl:greeks}{vega} du premier mois vaut 5.064, alors que le bump de la volatilité de ce mois avec
les mêmes tirages donne 5.065, et les \omterm{def:dv:greeks-and-the-hedging-pnl:greeks}{vegas} tombent à 0.016 pour le dernier mois (\cref{fig:dv:monte-carlo-pricers-in-practice:adjoint}).
\end{example}

\begin{omfigure}
\begin{tikzpicture}
\begin{axis}[width=10.5cm,height=4.6cm,ybar,bar width=12pt,
  xlabel={compartiment de volatilité (mois)},ylabel={vega (par unité de volatilité)},
  tick label style={font=\small},label style={font=\small},
  xmin=0.3,xmax=12.7,ymin=0,ymax=5.6,grid=major,grid style={gray!25},xtick={1,2,3,4,5,6,7,8,9,10,11,12},
  table/col sep=comma]
\addplot[omThm,fill=omThm!40] table[x=month,y=vega]{figdata/derivatives/23-monte-carlo-pricers-in-practice/bucket_vega.csv};
\end{axis}
\end{tikzpicture}

\omcaption{\omterm{def:dv:greeks-and-the-hedging-pnl:greeks}{Vegas} par compartiment d'une \omterm{def:dv:asians-lookbacks-cliquets-and-forward-starts:asian}{option asiatique} arithmétique d'un an, obtenus par une passe adjointe le long de chaque trajectoire~: la sensibilité de la moyenne à la volatilité précoce
est la plus forte, car les mouvements précoces se répercutent sur chaque constatation ultérieure. Données~: le tutoriel.}\label{fig:dv:monte-carlo-pricers-in-practice:adjoint}
\end{omfigure}

\section{Bridges et nombres quasi-aléatoires}

Les suites à discrépance faible (livre~4, chapitre~26) remplissent le cube unité plus uniformément que les nombres aléatoires, et, pour des intégrandes régulières, leur erreur décroît plus vite
que $1/\sqrt N$. Une trajectoire de $d$ dates est toutefois un point en $d$ dimensions, et les suites à discrépance faible perdent leur avantage dans les dernières
coordonnées. Le pont brownien réordonne la construction de sorte que les premières coordonnées portent l'essentiel de la variance de la trajectoire. La première normale fixe
l'extrémité $W_T$, la seconde le milieu conditionnellement aux extrémités, et ainsi de suite. Les bonnes premières coordonnées de la suite déterminent alors la forme de la trajectoire.

\begin{example}[Une option asiatique avec des points de Halton]\label{ex:dv:monte-carlo-pricers-in-practice:qmc}
Un call asiatique à 16 dates (un an, strike 100, 20~\%). Avec 4\,096 trajectoires, l'erreur type Monte-Carlo vaut 0.131. Les points de Halton, décalés aléatoirement vingt fois
pour mesurer leur erreur, donnent 0.027 avec la construction usuelle pas à pas et 0.0092 avec le pont brownien, quatorze fois plus faible que Monte-Carlo,
soit la précision d'environ 200 fois plus de trajectoires aléatoires. Avec 256 trajectoires, les trois erreurs valent 0.53, 0.34 et 0.13
(\cref{fig:dv:monte-carlo-pricers-in-practice:qmc}).
\end{example}

\begin{omfigure}
\begin{tikzpicture}
\begin{axis}[width=10.5cm,height=5cm,xmode=log,ymode=log,log basis x=2,
  xlabel={trajectoires},ylabel={erreur (écart-type)},
  tick label style={font=\small},label style={font=\small},
  xmin=200,xmax=5200,ymin=0.005,ymax=0.8,grid=major,grid style={gray!25},
  xtick={256,1024,4096},xticklabels={256,1\,024,4\,096},
  legend columns=3,legend cell align=left,legend style={font=\footnotesize,at={(0.5,-0.3)},anchor=north,draw=none,fill=none,/tikz/every even column/.append style={column sep=8pt}},
  table/col sep=comma]
\addplot[omProp,very thick,mark=*,mark size=1.6pt] table[x=n,y=mc]{figdata/derivatives/23-monte-carlo-pricers-in-practice/qmc.csv};
\addplot[omDef,very thick,mark=square*,mark size=1.6pt] table[x=n,y=halton]{figdata/derivatives/23-monte-carlo-pricers-in-practice/qmc.csv};
\addplot[omThm,very thick,mark=triangle*,mark size=1.8pt] table[x=n,y=bridge]{figdata/derivatives/23-monte-carlo-pricers-in-practice/qmc.csv};
\legend{Monte Carlo,Halton,Halton avec pont brownien}
\end{axis}
\end{tikzpicture}

\omcaption{Erreur d'un call asiatique à 16 dates en fonction du nombre de trajectoires~: erreur type de Monte-Carlo, et dispersion sur les décalages aléatoires des points de Halton
selon la construction pas à pas et la construction par pont brownien. Données~: le tutoriel.}\label{fig:dv:monte-carlo-pricers-in-practice:qmc}
\end{omfigure}

\section{Concevoir un pricer Monte-Carlo}

Un moteur de production sépare ce qui varie de ce qui est partagé. Les générateurs de trajectoires sont un par modèle~: lognormal, \omterm{def:dv:local-volatility:model}{volatilité locale}, Heston-QE,
stochastic-local. Ils écrivent les trajectoires sur une grille de dates dont le produit a besoin, et ils ne connaissent jamais le paiement. Les paiements ne lisent que les trajectoires (chapitres 16 et
18). Les estimateurs, la \omterm{def:dv:monte-carlo-pricers-in-practice:lsm}{méthode de Longstaff--Schwartz}, les \omterm{def:dv:greeks-and-the-hedging-pnl:greeks}{grecques} et la \omterm{def:dv:monte-carlo-pricers-in-practice:dual}{borne supérieure duale} se placent au-dessus. Quatre règles le gardent rigoureux.
\begin{itemize}
\item \textbf{Les graines sont des données.} Tout prix est reproductible à partir de sa graine. Les \omterm{def:dv:greeks-and-the-hedging-pnl:greeks}{grecques} par bump utilisent les mêmes tirages (nombres aléatoires communs).
\item \textbf{Tout nombre porte son erreur.} L'erreur type pour les statistiques, un contrôle de convergence pour la discrétisation, et un écart de bornes pour
l'exercice anticipé.
\item \textbf{Réduire la variance d'abord.} Variables antithétiques, variables de contrôle (l'\omterm{def:dv:asians-lookbacks-cliquets-and-forward-starts:asian}{option asiatique} géométrique du chapitre~16), conditionnement par la \omterm{prop:dv:stochastic-volatility:mixing}{formule de mélange} du chapitre~10 et
le bridge rapporte en général davantage que des trajectoires supplémentaires.
\item \textbf{Vectoriser sur les trajectoires, boucler sur le temps.} Les générateurs de la construction sont des tableaux de trajectoires avancés date par date, la forme naturelle pour numpy, pour
un GPU, et pour la passe adjointe qui remonte sur les mêmes dates.
\end{itemize}

\section{Tutoriel~: bornes, grecques et bridges}\label{tut:dv:monte-carlo-pricers-in-practice:tut}

\begin{tutorial}
\textbf{Objectif.} Pricer un max-call bermudéen par Longstaff--Schwartz et le borner par le haut à l'aide de la borne duale~; calculer par les \omterm{def:dv:monte-carlo-pricers-in-practice:pathwise}{grecques trajectorielles} et par rapport de vraisemblance les \omterm{def:dv:greeks-and-the-hedging-pnl:greeks}{grecques} d'un call et d'un digital, ainsi que les \omterm{def:dv:greeks-and-the-hedging-pnl:greeks}{vegas} par compartiment au moyen d'une passe adjointe~;
comparer Euler et QE pour Heston~; mesurer l'effet du pont brownien sur les
points de Halton. \textbf{État final~:} les quatre figures, le tableau et les valeurs numériques du problème du week-end.
\begin{tutsteps}
\item \textbf{La régression}, à rebours sur les dates d'exercice, sur les trajectoires dans la monnaie uniquement~:
\omcode{code/firm/mcpricer/firm_mcpricer.py}{85}{103}{Régression de Longstaff--Schwartz.}\label{lst:dv:monte-carlo-pricers-in-practice:lsm}
\item \textbf{La borne duale}~: valeurs de continuation internes, accroissements de martingale $L_j-Q_{j-1}$, le maximum~:
\omcode{code/firm/mcpricer/firm_mcpricer.py}{126}{157}{La borne supérieure d'Andersen--Broadie.}\label{lst:dv:monte-carlo-pricers-in-practice:dual}
\item \textbf{La passe adjointe} pour les \omterm{def:dv:greeks-and-the-hedging-pnl:greeks}{vegas} par compartiment~:
\omcode{code/firm/mcpricer/firm_mcpricer.py}{184}{206}{Passes directe et rétrograde pour une option d'achat asiatique.}\label{lst:dv:monte-carlo-pricers-in-practice:adjoint}
\item \textbf{Lancer} \texttt{dv\_mc.bounds()}, \texttt{put\_bias()}, \texttt{greek\_table()}, \texttt{adjoint()}, \texttt{heston\_bias()}, \texttt{qmc\_study()}
et \texttt{fig\_mc.py}.
\end{tutsteps}
\textbf{Pour aller plus loin.} Pricer le max-call avec une corrélation de 0.5~; ajouter l'\omterm{def:dv:asians-lookbacks-cliquets-and-forward-starts:asian}{option asiatique} géométrique comme variable de contrôle en plus des points de Halton~; remplacer le
rapport de vraisemblance du digital par une espérance conditionnelle sur le dernier pas de temps (lissage par conditionnement).
\end{tutorial}

\section{Construction~: le moteur de Monte-Carlo}\label{bld:dv:monte-carlo-pricers-in-practice:mcpricer}

\begin{build}
\bfield{Rôle} Le moteur de simulation de la firme miniature~: générateurs de trajectoires, exercice anticipé avec bornes, \omterm{def:dv:greeks-and-the-hedging-pnl:greeks}{grecques} par trois méthodes, nombres quasi-aléatoires~;
le moteur de Monte-Carlo du chapitre~28 l'enveloppe.
\bfield{Interface} \texttt{gbm\_paths(s0, vols, corr, times, r, q, n, seed)}~; \texttt{heston\_terminal(\ldots, scheme)}~; \texttt{lsm\_fit(paths, ex\_idx, payoff,
basis, k, disc)}, \texttt{lsm\_price}, \texttt{exercise\_now}, \texttt{max\_call\_basis}, \texttt{small\_basis}~; \texttt{dual\_upper(outer, \ldots, step\_sim,
n\_inner, seed)}~; \texttt{call\_greeks(\ldots, digital, smooth)}~; \texttt{asian\_adjoint}~; \texttt{halton(n, dim, shift)}, \texttt{brownian\_bridge(z, t)},
\texttt{norm\_inv}.
\bfield{Règles} Régression et évaluation sur des trajectoires indépendantes~; une borne inférieure est publiée avec sa \omterm{def:dv:monte-carlo-pricers-in-practice:dual}{borne supérieure duale} pour chaque option bermudéenne~; \omterm{def:dv:monte-carlo-pricers-in-practice:pathwise}{grecques trajectorielles} là où
le paiement est continu, par rapport de vraisemblance ou lissage là où il ne l'est pas~; graines explicites~; aucune constante mémorisée dans le calcul numérique (l'inverse de la loi normale est
la méthode de Newton appliquée à la fonction d'erreur).
\bfield{Tests d'acceptation} \texttt{code/firm/mcpricer/tests/}~: les bornes bermudéennes encadrent la valeur de l'arbre et la borne duale diminue avec le nombre de trajectoires internes~; \omterm{def:dv:monte-carlo-pricers-in-practice:pathwise}{grecques trajectorielles} et
\omterm{def:dv:monte-carlo-pricers-in-practice:lr}{grecques par rapport de vraisemblance} s'accordent avec Black--Scholes et le digital trajectoriel est nul~; la passe adjointe coïncide avec le bump~; QE coïncide avec le prix de Fourier~; les premiers
points de Halton, l'inverse de la loi normale et les covariances du pont brownien.
\bfield{Pour aller plus loin} Points de Sobol avec des nombres de direction publiés (cités, non rappelés)~; Monte-Carlo multiniveau (livre~4)~; une passe adjointe à travers
Longstaff--Schwartz.
\end{build}

\begin{omsources}
\item F.~A.~Longstaff and E.~S.~Schwartz, «~Valuing American options by simulation~: a simple least-squares approach~», \emph{Review of Financial Studies}
14(1) (2001) 113--147.
\item L.~Andersen and M.~Broadie, «~Primal-dual simulation algorithm for pricing multidimensional American options~», \emph{Management Science} 50(9)
(2004) 1222--1234.
\item L.~Andersen, «~Simple and efficient simulation of the Heston stochastic volatility model~», \emph{Journal of Computational Finance} 11(3) (2008) 1--42.
\item M.~Broadie and P.~Glasserman, «~Estimating security price derivatives using simulation~», \emph{Management Science} 42(2) (1996) 269--285.
\item P.~Glasserman, \emph{Monte Carlo Methods in Financial Engineering} (Springer, 2003).
\item M.~Giles and P.~Glasserman, «~Smoking adjoints~: fast Monte Carlo Greeks~», \emph{Risk} 19 (2006) 88--92~; rapport Oxford NA-05-15 (2005).
\end{omsources}

\section{Exercices}

\begin{exercise}[$\star$]\label{exo:dv:monte-carlo-pricers-in-practice:1}
Expliquer pourquoi un prix de Longstaff--Schwartz, calculé sur des trajectoires indépendantes de la régression, constitue une borne inférieure.
\end{exercise}

\begin{exercise}[$\star$]\label{exo:dv:monte-carlo-pricers-in-practice:2}
Établir le delta trajectoriel d'un call sous Black--Scholes, $e^{-rT}\mathbf 1\{S_T>K\}S_T/S_0$.
\end{exercise}

\begin{exercise}[$\star$]\label{exo:dv:monte-carlo-pricers-in-practice:3}
Expliquer pourquoi le delta trajectoriel d'un digital est nul, et pourquoi le delta par rapport de vraisemblance ne l'est pas.
\end{exercise}

\begin{exercise}[$\star\star$]\label{exo:dv:monte-carlo-pricers-in-practice:4}
Établir le poids du \omterm{def:dv:greeks-and-the-hedging-pnl:greeks}{vega} par rapport de vraisemblance $(Z^2-1)/\sigma-Z\sqrt T$ pour un prix terminal lognormal.
\end{exercise}

\begin{exercise}[$\star\star$]\label{exo:dv:monte-carlo-pricers-in-practice:5}
Expliquer pourquoi un nombre insuffisant de trajectoires internes biaise la borne duale vers le haut, jamais vers le bas.
\end{exercise}

\begin{exercise}[$\star\star$]\label{exo:dv:monte-carlo-pricers-in-practice:6}
Expliquer pourquoi le pont brownien aide les nombres quasi-aléatoires mais non les nombres pseudo-aléatoires.
\end{exercise}

\begin{exercise}[$\star\star\star$]\label{exo:dv:monte-carlo-pricers-in-practice:7}
\emph{Programmation.} Ajuster et évaluer la politique d'exercice de Longstaff--Schwartz du max-call sur les mêmes trajectoires, et comparer avec le prix sur trajectoires indépendantes. Expliquer le signe de la
différence.
\end{exercise}

\begin{exercise}[$\star\star\star$]\label{exo:dv:monte-carlo-pricers-in-practice:8}
\emph{Trouver le défaut.} «~Notre Monte-Carlo évalue l'option bermudéenne à deux centimes près~: l'erreur type vaut 0.01.~»
\end{exercise}

\section{Problème~: les bornes bermudéennes}

\begin{problem}[{Problème du week-end --- qualité de la politique d'exercice}]\label{pb:dv:monte-carlo-pricers-in-practice:1}
Un desk vend un call bermudéen de trois ans sur le meilleur de trois actions (chacune à 100, strike 100, volatilité 20~\%, rendement de dividende 10~\%, taux 5~\%, neuf dates d'exercice)
et le valorise par la \omterm{def:dv:monte-carlo-pricers-in-practice:lsm}{méthode de Longstaff--Schwartz}.

\medskip\noindent\textbf{Partie I --- La borne inférieure.}
\begin{enumerate}
\item Donner le prix avec la base riche et son erreur type.
\item Expliquer pourquoi l'ajustement et l'évaluation se font sur des trajectoires distinctes.
\item Donner le prix avec la base réduite (le plus grand prix et son carré).
\item Identifier les décisions d'exercice que la base réduite traite mal.
\item Préciser ce que révèle le prix d'un max-call européen.
\end{enumerate}

\medskip\noindent\textbf{Partie II --- La borne supérieure.}
\begin{enumerate}[resume]
\item Donner la borne duale pour la base riche avec 1\,000 trajectoires externes et 1\,000 trajectoires internes, ainsi que son erreur type.
\item Donner l'écart et l'exprimer en proportion du prix.
\item Donner la borne et l'écart pour la base réduite.
\item Pour le put bermudéen à un seul actif, expliquer comment la borne varie avec 100, 400 et 1\,600 trajectoires internes.
\item Expliquer ce que coûte la borne duale par rapport à la borne inférieure, et pourquoi.
\end{enumerate}

\medskip\noindent\textbf{Partie III --- Autour du prix.}
\begin{enumerate}[resume]
\item Expliquer comment calculer le delta de l'option bermudéenne.
\item Expliquer pourquoi éviter le \omterm{def:dv:greeks-and-the-hedging-pnl:bump}{bump-and-reprice} avec de nouvelles trajectoires.
\item Préciser ce qu'apporte la méthode adjointe à un portefeuille de telles transactions.
\item Expliquer en quoi les nombres quasi-aléatoires aideraient, et ce qui les limite ici.
\item Identifier l'erreur de discrétisation que ce produit n'a pas, et le modèle qui en introduirait une.
\end{enumerate}

\medskip\noindent\textbf{Partie IV --- Jugement.}
\begin{enumerate}[resume]
\item Déterminer quel prix coter et lequel passer en compte~?
\item Préciser la réserve à constituer pour la politique d'exercice~?
\item Préciser à partir de quel moment arrêter d'améliorer la base~?
\item Énoncer le \emph{résultat nommé}~: l'écart entre la borne inférieure de Longstaff--Schwartz et la \omterm{def:dv:monte-carlo-pricers-in-practice:dual}{borne supérieure duale} pour le max-call bermudéen à trois actifs, et
sa sensibilité à la base de régression.
\item En une phrase~: que mesure l'écart entre les bornes~?
\end{enumerate}
\end{problem}

\section{Questions d'entretien}

\begin{interviewq}[$\star$ \iqroles{développeur, chercheur}]\label{iq:dv:monte-carlo-pricers-in-practice:1}
Expliquez l'algorithme de Longstaff--Schwartz.
\end{interviewq}

\begin{interviewq}[$\star\star$ \iqroles{chercheur}]\label{iq:dv:monte-carlo-pricers-in-practice:2}
Comment obtenez-vous une borne supérieure pour un prix bermudéen par simulation~?
\end{interviewq}

\begin{interviewq}[$\star\star$ \iqroles{chercheur}]\label{iq:dv:monte-carlo-pricers-in-practice:3}
Comparez les \omterm{def:dv:monte-carlo-pricers-in-practice:pathwise}{grecques trajectorielles} et les \omterm{def:dv:monte-carlo-pricers-in-practice:lr}{grecques par rapport de vraisemblance}. Quand chacune échoue-t-elle~?
\end{interviewq}

\begin{interviewq}[$\star\star$ \iqroles{développeur}]\label{iq:dv:monte-carlo-pricers-in-practice:4}
Comment simuleriez-vous efficacement le \omterm{def:dv:stochastic-volatility:heston}{modèle de Heston}, et pourquoi pas Euler~?
\end{interviewq}

\begin{interviewq}[$\star\star$ \iqroles{développeur, risque}]\label{iq:dv:monte-carlo-pricers-in-practice:5}
Votre risque Monte-Carlo est bruité d'un jour à l'autre. Que changez-vous~?
\end{interviewq}

\begin{interviewq}[$\star\star\star$ \iqroles{développeur}]\label{iq:dv:monte-carlo-pricers-in-practice:6}
Expliquez la différentiation algorithmique adjointe pour les \omterm{def:dv:greeks-and-the-hedging-pnl:greeks}{grecques} Monte-Carlo. Quel en est le coût en mémoire et en temps~?
\end{interviewq}
